\begin{abstract}
Laser-induced breakdown spectroscopy (LIBS) can in principle quantify elemental
composition without matrix-matched standards by inverting a physical model of
the plasma emission, the calibration-free (CF-LIBS) programme. In practice the
inversion is ill-posed: temperature, electron density, optical depth, detector
response and line identification are entangled, and small errors in any one
appear as large, self-consistent errors in composition. This paper sets out the
theoretical basis, the explicit assumptions, the processing workflow, the
datasets and the fit metrics that underpin \alibz, an open-source toolkit that
treats a LIBS spectrum as a non-negative linear combination of forward-modelled
element--ion templates evaluated at a shared plasma state, and solves the
inverse problem with physical constraints rather than empirical calibration
curves. We describe the local-thermodynamic-equilibrium emission model and its
conventions, the treatment of self-absorption by resolved-profile and
doublet-ratio inversion, the electron-density anchors, the stage-consistency
thermometer that fixes temperature where the amplitude data cannot, the
peak-extraction and refinement pipeline, the whole-pattern weighted
non-negative solve with its evidence terms and basin guards, and the
true-negative contest that resolves spectral confounders. We define the
diagnostics by which a fit is judged, report where the method currently
succeeds and fails on synthetic and real spectra from a handheld analyser, and
state the validation gates that remain open: composition invariance between
argon and air, pure-element standards for the plasma physics, and certified
reference materials for the composition solve.
\end{abstract}

\noindent\textbf{Keywords:} laser-induced breakdown spectroscopy; calibration-free
LIBS; Saha--Boltzmann; self-absorption; non-negative least squares; spectral
unmixing; handheld LIBS
